% ==============================================================
\section{結論}\label{sec:conclusion}
% ==============================================================

本研究では\emph{稀少密集パターン}（RDP）マイニング問題を導入した。これは二つの条件を同時に満たすアイテムセットを対象とする新しい定式化である：大域的稀少性（全体サポートが低い）と局所的時間密集性（少なくとも一つのスライディングウィンドウ内で共起カウントが高い）。
このパターンクラスは、プロモーションバースト、セキュリティインシデント、季節的同時購入など、標準的な頻出パターンマイニング（稀少パターンを枝刈りしてしまう）と標準的な稀少パターンマイニング（時間的構造を無視する）の両方からは見えない、実用上重要な現象を捉える。

本研究の理論的な主要貢献は\emph{弱反単調性}である：RDP条件の組み合わせは単調でも反単調でもないが、局所密集性成分のみは反単調である。
これにより、まず局所的に密集するすべてのパターンを発見し、次に大域的稀少性でフィルタリングする\emph{二段階マイニング}アルゴリズムにおけるApriori型枝刈りが可能となる。
本アルゴリズムが健全かつ完全であること---すべてのRDPを過不足なく発見すること---を証明した。

合成データ上の実験結果は、正解RDPの100\%回収、稀少性フィルタによる80\%以上の偽陽性削減、10万トランザクションまでのほぼ線形なスケーラビリティを実証した。
さらに、実小売データ上での評価により、季節的バンドルやプロモーション関連パターンなど、既存手法では発見困難な実用的パターンの検出を確認した。

\paragraph{制約}
現在のアルゴリズムは固定のウィンドウサイズ $W$ と密集閾値 $\theta$ を前提とする。
データ特性に基づくこれらのパラメータの自動選択は未解決の課題である。
大域的稀少性条件は単一の閾値 $\sigma_{\max}$ を使用するが、アイテムごとまたはアイテムセットサイズごとの閾値により柔軟性を向上できる可能性がある。

\paragraph{今後の研究}
以下の方向性について検討の価値がある：
\begin{enumerate}
\item \textbf{マルチスケール分析}：複数のウィンドウサイズにわたるRDPの同時検出。Kleinbergの多レベルバースト検出~\cite{kleinberg2003bursty}と類似の階層的分解の利用が考えられる。

\item \textbf{ストリーミング設定}：トランザクションが連続的に到着するオンラインマイニングへの適応。RDPをインクリメンタルに維持する。特にネットワーク侵入検知ではリアルタイム性が不可欠であり、スライディングウィンドウのインクリメンタル更新と稀少性フィルタの遅延評価を組み合わせたオンラインRDP検出が重要な課題となる。

\item \textbf{稀少性を考慮した枝刈り}：第1段階において稀少性条件を活用するより厳密な枝刈り戦略の開発。大域サポートの上界推定を早期に行い、$\sigma_{\max}$を超えるパターンを第1段階で枝刈りすることで、候補集合を削減できる可能性がある。

\item \textbf{イベント帰属}：発見されたRDPを外部イベント（プロモーション、インシデント等）と関連付け、稀少パターンが特定の時点で局所的に密集した\emph{理由}を説明する。

\item \textbf{ネットワーク侵入検知への適用}：ネットワークログにおいて、攻撃パターン（ポートスキャン、DDoS等）は大域的に稀少だが攻撃期間中に局所的に密集する。NSL-KDDやCICIDS2017等の公開侵入検知データセットにおけるRDPの有効性検証が具体的な次のステップである。

\item \textbf{医療データ分析}：電子カルテにおける稀少な副作用共起パターンのバースト検出や、新薬承認後の処方パターン変化の追跡。MIMIC-IVデータベースを用いた検証が計画されている。プライバシー保護のための差分プライバシーとの統合も重要な課題である。

\item \textbf{統計的有意性検定の統合}：現在の評価は5シードの平均と標準偏差に基づくが、パターンレベルでの有意性検定（Westfall--Young法~\cite{szathmary2007rare}等）の統合により、偽発見率を制御した上でのRDP報告が可能となる。
\end{enumerate}
